\scanpage{74}
% archive sheet number: 73

\unsure{fonctoriellt|fonctoriellement} de $X$, mais on trouverait encore un iso fonctoriel en $X$, $\mathcal F_X\simeq\underline\omega^{\otimes2}_{X/\mathbb C}$. \deleted{De} \unsure{$=$|''} dans (\unsure{III|II}). Ainsi les 2-formes sont essentielles.

2) \deleted{G} Le torseur $\mathcal T$ dont il est question dans II (et en marge de I) dépend du choix de $G$ (et de $\mathcal F$, si on veut -- c'est alors un torseur sous $\mathcal F$) mais les \inserted{torseurs associés aux} différents choix sont can. (semi)-isomorphes, de façon précise

\[
\mathcal T^G\xrightarrow{\sim}\mathcal T^{\varphi G}
\]

\unsure{corresp. avec|comp. avec|au-dessus de}

\[
\underline\omega_X^{\otimes2}\xrightarrow{\varphi.\mathrm{id}}\underline\omega_X^{\otimes2}
\]

i.e. $\mathcal T^{\varphi G}$ est ``multiplié'' de $\unsure{\mathcal E^G|\mathcal T^G}$ par $\varphi$. Ainsi, le choix ``le plus fin'' de $\mathcal T$ est $\mathcal T^{G_0}$ (où $G_0$ est défini en (III)), alors le $\mathcal T^G$ correspondant \deleted{au choix} \inserted{$G$} de (\unsure{II|III}) est un multiple 12 ième de $\mathcal T^{G_0}$.

3) On \unsure{peut|pourrait} se \deleted{demander} \inserted{poser} la question de caractériser $G_0(f,g)$ \uline{en une caractéristique donnée} $p>0$. On peut prendre les conditions a) \inserted{\unsure{reprises|reçues} comme plus haut} b) c), plus d') que $G_0(f,g)$ n'est pas id. nul -- \uline{cela caractérise} $G_0(f,g)$ à multiplication près par une constante $c\in\mathbb F_p^*$. \unsure{M̄|M.} précis$^{\mathrm{t}}$, si on se fixe un schéma de base absolu \uline{intègre} $S_0$ (pour les $S$ et les $X$), l'indétermination est dans $\Gamma(S,\mathcal O_S^*)$. Si $S$ n'est pas intègre, il faut supposer que, sur chaque $k(s)$, l'op. induit (pour des $X$ lisses \deleted{\unsure{sur le|sur une}} \inserted{de dim$^{\mathrm{ion}}$} rel. 1 sur $k(s)$, )

% slanted
\margin{\uline{faux} ?! \unsure{car un|par un|si} \unsure{$D_W$|sur $W$} ou \unsure{déjà|déf.|diff.} \unsure{$G$} d'ordre $\leq3$)}

% margin note written sideways
\margin{\unsure{$d$|$G$|d')} n'est pas identiquement nul.}
